\chapter{บทนำ}
\label{Ch:Introduction}
ในปัจจุบัน ปัญหาการเปลี่ยนแปลงภูมิอากาศโลก เป็นประเด็นสำคัญที่ส่งผลกระทบทั้งต่อระบบนิเวศของสิ่งมีชีวิตและต่อระบบเศรษฐกิจ ซึ่งมีผลโดยตรงต่อคุณภาพชีวิตของผู้คนทั่วโลก ประเด็นดังกล่าวได้ก่อให้เกิดความตระหนักรู้และการร่วมมือกันในการปกป้อง พิทักษ์ และรักษาทรัพยากรธรรมชาติให้ยั่งยืน ในภาคส่วนซอฟต์แวร์และเทคโนโลยีสารสนเทศก็มีบทบาทสำคัญ เนื่องจากการพัฒนาซอฟต์แวร์ที่ไม่มีประสิทธิภาพอาจนำไปสู่การใช้พลังงานที่สูงเกินความจำเป็น ดังนั้นผู้จัดทำจึงมุ่งเน้นที่จะศึกษาและพัฒนา แนวทางการเขียนโปรแกรมเชิง Green Coding ผ่านภาษาคอมพิวเตอร์ที่ได้รับความนิยมอย่าง Python เพื่อเสนอเครื่อง
มือที่สามารถช่วยวิเคราะห์ผลกระทบต่อสิ่งแวดล้อมได้

ในบทนี้จะอธิบายถึงที่มาและความสำคัญของปัญหา วัตถุประสงค์ ขอบเขตการดำเนินงาน และประโยชน์ที่คาดว่าจะได้รับจากการศึกษาและพัฒนาเครื่องมือนี้

% -----------------------------------------------------

\section{ที่มาและความสำคัญของปัญหา}
\label{Sec:ProblemStatement}
ในปัจจุบันอุตสาหกรรมซอฟต์แวร์มีอัตราการเติบโตอย่างต่อเนื่อง ส่งผลให้ความต้องการใช้พลังงานสำหรับการประมวลผลเพิ่มสูงขึ้นตามไปด้วย การใช้พลังงานดังกล่าวยังส่งผลต่อทรัพยากรที่มีอยู่จำกัด ส่งผลกระทบต่อสิ่งแวดล้อมและการเปลี่ยนแปลงสภาพภูมิอากาศด้วยเช่นกัน

ผู้จัดทำเล็งเห็นถึงปัญหาการใช้พลังงานที่มากกว่าที่ควรซึ่งเกิดจากการพัฒนาซอฟต์แวร์ที่ไม่คำนึงถึงผลกระทบต่อการใช้พลังงาน จึงมุ่งเน้นการศึกษาการทำงานของ Green Code Smell เพื่อเป็นเครื่องมือช่วยตรวจสอบการทำงานของโค้ดในภาษา Python โดยจัดทำผลลัพธ์ในรูปแบบ Console Log, JSON และการแสดงผลเชิงภาพ (Visualization) ที่จะสะท้อนถึงประสิทธิภาพด้านการใช้พลังพลังงานและผลกระทบต่อสิ่งแวดล้อม อันจะนำไปสู่การส่งเสริมการพัฒนาซอฟต์แวร์ที่เป็นมิตรต่อสิ่งแวดล้อมมากยิ่งขึ้น

% -----------------------------------------------------

\section{วัตถุประสงค์}
\label{Sec:Objective}
โครงการนี้มีวัตถุประสงค์เพื่อสนับสนุนการพัฒนาซอฟต์แวร์อย่างยั่งยืน โดยมุ่งเน้นการตรวจจับและทำความเข้าใจรูปแบบโค้ดที่ส่งผลต่อการใช้พลังงาน (Green Code Smells) ในภาษา Python ผ่านการดำเนินการดังต่อไปนี้
\begin{enumerate}
    \item {\boldfont พัฒนาไลบรารีภาษา Python และส่วนขยายสำหรับ Visual Studio Code} \\
    เพื่อใช้ในการตรวจจับ Green Code Smells และรูปแบบการเขียนโปรแกรมอื่น ๆ ที่อาจส่งผลกระทบอย่างมีนัยสำคัญต่อการใช้พลังงานของซอฟต์แวร์ โดยรองรับการวิเคราะห์โค้ดอัตโนมัติและการประเมินเชิงโครงสร้าง
    \item {\boldfont ศึกษาผลกระทบด้านพลังงานของ Green Code Smells และพัฒนาวิธีการตรวจจับที่มีประสิทธิภาพ} \\
    โดยมุ่งเน้นการวิเคราะห์และทำความเข้าใจความสัมพันธ์ระหว่างรูปแบบโค้ดกับปริมาณการใช้พลังงาน รวมถึงการพัฒนา หรือการตรวจสอบความถูกต้องของวิธีการวิเคราะห์ที่ใช้ในการตรวจจับรูปแบบดังกล่าว
    \item {\boldfont สรุปผลและจัดทำข้อเสนอแนะเพื่อสนับสนุนการพัฒนาซอฟต์แวร์ที่ยั่งยืน} \\
    ผ่านการจัดทำรายงาน การแสดงผลผลการวิเคราะห์ และการประยุกต์ใช้ตัวชี้วัดด้านความยั่งยืนและการใช้พลังงาน เพื่อสนับสนุนให้ผู้พัฒนาสามารถปรับปรุงคุณภาพโค้ดตามแนวทาง Green Coding ได้อย่างมีประสิทธิภาพ
\end{enumerate}

% -----------------------------------------------------

\section{ขอบเขตงาน}
\label{Sec:ScopeOfWork}
โครงการนี้กำหนดขอบเขตการดำเนินงานเพื่อรองรับการพัฒนาและประเมินผล Green Code Smell Detector สำหรับภาษา Python โดยครอบคลุมประเด็นสำคัญดังต่อไปนี้
\begin{enumerate}
    \item {\boldfont การพัฒนาไลบรารีภาษา Python และส่วนขยายสำหรับ Visual Studio Code} \\
    พัฒนาไลบรารีและส่วนขยายเพื่อรองรับการตรวจสอบ Green Code Smells โดยใช้งานในเครื่องมือสำหรับการพัฒนาซอฟต์แวร์ที่มีความเป็นสากลอย่างเช่น Visual Studio Code
    \item {\boldfont รองรับการวิเคราะห์โค้ดด้วยโครงสร้าง Abstract Syntax Tree (AST)} \\
    ไลบรารีต้องสามารถประมวลผลและแยกวิเคราะห์โค้ดผ่านโครงสร้าง AST เพื่อสกัดองค์ประกอบของโค้ด เช่น Node, Statement และ Structure ต่าง ๆ เพื่อใช้ในการตรวจจับรูปแบบโค้ดที่ไม่ประหยัดพลังงาน
    \item {\boldfont ชุดของกฎเกณฑ์และการตรวจจับ Green Code Smell} \\
    กำหนดชุดของกฎเกณฑ์ (Ruleset) ตรรกะการสำหรับการวิเคราะห์ และเงื่อนไขเบื้องต้นสำหรับการตรวจจับ Green Code Smells โดยให้สอดคล้องกับแนวคิดเชิงโครงสร้างของภาษา Python และข้อมูลจากงานวิจัยที่เกี่ยวข้อง
    \item {\boldfont การประเมินและสรุปผลการวิเคราะห์โค้ด} \\
    ทำการสรุปผลการตรวจสอบโค้ด พร้อมจัดทำรายงานที่แสดงตัวชี้วัดที่สำคัญ เช่น การปล่อยคาร์บอน (Carbon Emission) การใช้พลังงาน (Energy Consumption) และตัวชี้วัดด้านความยั่งยืนอื่น ๆ เพื่อสนับสนุนการประเมินคุณภาพของโค้ด
\end{enumerate}

% -----------------------------------------------------

\section{ประโยชน์ที่คาดว่าจะได้รับ}
\label{Sec:ExpectedBenefit}
จากการดำเนินโครงการ มีความคาดหวังว่าจะได้ไลบรารี ที่สามารถใช้ตรวจสอบและวิเคราะห์ Green Code Smell ในภาษา Python ได้อย่างเป็นระบบ ซึ่งจะช่วยให้ผู้พัฒนาซอฟต์แวร์มีความตระหนักถึงผลกระทบต่อพลังงานและสิ่งแวดล้อมจากการเขียนโปรแกรม นอกจากนี้ยังสามารถใช้เป็นแนวทางเริ่มต้นในการประยุกต์แนวคิด Green Coding กับซอฟต์แวร์ประเภทอื่นในอนาคต โดยส่งเสริมการลดการใช้พลังงานที่ไม่จำเป็นและสนับสนุนการพัฒนาซอฟต์แวร์ที่ยั่งยืนและเป็นมิตรต่อสิ่งแวดล้อม

% -----------------------------------------------------
